\documentclass[a4paper]{article}
\usepackage[a4paper,top=15mm,bottom=15mm,left=20mm,right=20mm]{geometry}
\usepackage{polyglossia}
\setdefaultlanguage[babelshorthands=true]{russian}
\setotherlanguage{english}
\usepackage{fontspec}
\setmainfont{FreeSerif}
\newfontfamily{\russianfonttt}[Scale=0.85]{DejaVuSansMono}
\usepackage[tiny,compact]{titlesec}
\usepackage{titling}
\setlength{\droptitle}{-1cm}
\pretitle{\begin{center}\begin{bfseries}\Large}
\posttitle{\par\end{bfseries}\end{center}}
\preauthor{\begin{center}\normalsize}
\postauthor{\par\end{center}\vspace{-1.5cm}}
\usepackage{hyperref}
\usepackage{bookmark}
\usepackage{booktabs}
\usepackage{float}

\title{Примитивы линейной алгебры: сравнительный анализ форматов хранения разреженных матриц}
\author{Семенский Роман Николаевич}
\date{}

\begin{document}
\maketitle
\begin{flushright}
Группа: \emph{25.Б42-мм} \\
Кафедра: \emph{системного программирования} \\
Научный руководитель: \emph{Григорьев Семён Вячеславович} \\
Номер семестра практики: \emph{2}
\end{flushright}

\section{Введение}
Работа выполнялась в рамках учебной практики и проекта по реализации примитивов разреженной линейной алгебры. Разреженная матрица содержит преимущественно нулевые элементы. В формате разреженного хранения записываются ненулевые элементы и их индексы.

В проекте создана собственная реализация формата COO (Coordinate List~--- список координат) на языке~C. Её производительность сравнивалась с CSparse~--- библиотекой для операций над разреженными матрицами, входящей в SuiteSparse.

Результаты сравнения и профилирования предназначены для определения участков собственной реализации, требующих дальнейшей доработки.

\section{Цель и задачи работы}
Цель работы~--- реализовать операции умножения для разреженных матриц в формате COO, сравнить их время выполнения с CSparse и определить узкие места собственной реализации.

Для достижения цели были поставлены следующие задачи:
\begin{enumerate}
\item реализовать умножение двух COO-матриц;
\item реализовать умножение COO-матрицы на разреженный вектор без его преобразования в плотный массив;
\item подготовить матрицы разного размера из SuiteSparse Matrix Collection;
\item измерить время двух операций для COO и CSparse по единой методике;
\item измерить затраты времени по фазам собственных алгоритмов.
\end{enumerate}

\section{Описание решения}
В формате COO каждый ненулевой элемент задаётся тройкой «индекс строки, индекс столбца, значение». Тройки хранятся в трёх одномерных массивах. Для упорядочивания элементов по индексам строк и столбцов используется функция \texttt{qsort}.

При умножении матриц для второго сомножителя строится индекс строк. Значения одной строки результата накапливаются во временном массиве по индексам столбцов, после чего ненулевые значения переносятся в COO-матрицу результата.

Разреженный вектор представляется как COO-матрица размера $n\times1$. Для каждого ненулевого элемента матрицы значение вектора ищется двоичным поиском среди его ненулевых элементов. Плотный вектор из $n$ значений не создаётся.

Для сравнения CSparse получает те же матрицы в формате CSC (Compressed Sparse Column~--- сжатое хранение по столбцам). Входные представления подготавливаются до замеров; измеряется только операция умножения.

\section{Эксперименты}
\subsection{Методика измерений}
Цель сравнительного эксперимента~--- измерить время одного вызова умножения для COO-реализации и CSparse. Использовались девять матриц из SuiteSparse Matrix Collection. Набор включает матрицы размером от $62\times62$ до $9877\times9877$. Для симметричных файлов Matrix Market при загрузке явно добавляются элементы зеркальной части. Поэтому NNZ (Number of Non-Zeros~--- число ненулевых элементов) в таблицах указано после загрузки.

Измерялись $A\times A$ и $A\times x$. Вектор $x$ имел тот же размер, что и число столбцов $A$, а доля его ненулевых элементов составляла 10\%, поэтому вектор оставался разреженным. Перед созданием векторов генератор псевдослучайных чисел инициализировался фиксированным значением 42. Перед замерами выполнялись три прогревочных вызова. Затем собирались 15 выборок; число операций в каждой увеличивалось, пока её длительность не достигала 100~мс.

Характеристики вычислительной системы:
\begin{itemize}
\item операционная система: Ubuntu~26.04;
\item процессор: 13th~Gen Intel Core i5-13420H;
\item оперативная память: 16~ГБ;
\item компилятор: GCC (GNU Compiler Collection)~15.2.0, флаг оптимизации \texttt{-O3}.
\end{itemize}
Время измерялось монотонным таймером \texttt{clock\_gettime} с \texttt{CLOCK\_MONOTONIC}.

\subsection{Результаты сравнения}
В табл.~\ref{tab:mm} приведено время одного вызова $A\times A$. Замедление равно отношению времени COO ко времени CSparse.
\begin{table}[H]
\centering\small
\caption{Время умножения матриц $A\times A$}
\label{tab:mm}
\begin{tabular}{lrrrrr}
\toprule
Матрица & Размер & NNZ & COO, мкс & CSparse, мкс & Замедление, $\times$\\
\midrule
dolphins & $62\times62$ & 318 & 4{,}643 & 2{,}497 & 1{,}86\\
lesmis & $77\times77$ & 508 & 8{,}744 & 6{,}499 & 1{,}35\\
polbooks & $105\times105$ & 882 & 40{,}961 & 25{,}800 & 1{,}59\\
football & $115\times115$ & 1226 & 74{,}015 & 53{,}911 & 1{,}37\\
celegansneural & $297\times297$ & 2345 & 216{,}916 & 67{,}057 & 3{,}23\\
netscience & $1589\times1589$ & 5484 & 1106{,}743 & 67{,}400 & 16{,}42\\
add20 & $2395\times2395$ & 17319 & 2882{,}296 & 1045{,}834 & 2{,}76\\
ca-GrQc & $5242\times5242$ & 28980 & 13094{,}578 & 1377{,}051 & 9{,}51\\
ca-HepTh & $9877\times9877$ & 51971 & 43480{,}173 & 3218{,}179 & 13{,}51\\
\bottomrule
\end{tabular}
\end{table}

Для $A\times A$ CSparse имела меньшее время на всех девяти матрицах; замедление COO составляло от 1{,}35 до 16{,}42 раза.

В табл.~\ref{tab:mv} приведены результаты $A\times x$. В обеих реализациях вектор оставался разреженным.
\begin{table}[H]
\centering\small
\caption{Время умножения матрицы на вектор $A\times x$}
\label{tab:mv}
\begin{tabular}{lrrrrr}
\toprule
Матрица & Размер & NNZ & COO, мкс & CSparse, мкс & Замедление, $\times$\\
\midrule
dolphins & $62\times62$ & 318 & 1{,}605 & 0{,}160 & 10{,}03\\
lesmis & $77\times77$ & 508 & 2{,}444 & 0{,}141 & 17{,}33\\
polbooks & $105\times105$ & 882 & 4{,}639 & 0{,}213 & 21{,}78\\
football & $115\times115$ & 1226 & 6{,}445 & 0{,}200 & 32{,}23\\
celegansneural & $297\times297$ & 2345 & 21{,}385 & 0{,}329 & 65{,}00\\
netscience & $1589\times1589$ & 5484 & 101{,}257 & 1{,}489 & 68{,}00\\
add20 & $2395\times2395$ & 17319 & 341{,}174 & 3{,}044 & 112{,}08\\
ca-GrQc & $5242\times5242$ & 28980 & 1175{,}114 & 8{,}317 & 141{,}29\\
ca-HepTh & $9877\times9877$ & 51971 & 2499{,}207 & 16{,}740 & 149{,}30\\
\bottomrule
\end{tabular}
\end{table}

Для $A\times x$ CSparse имела меньшее время на всех девяти матрицах; замедление COO составляло от 10{,}03 до 149{,}30 раза.

\subsection{Профилирование COO-реализации}
Цель профилирования~--- определить фазы, которые фактически занимают наибольшую долю времени. Измерялись проверка упорядоченности входных данных, построение индекса, накопление произведений, просмотр буфера и создание результата. Профильная сборка использовала \texttt{-O3}.

Для профилирования $A\times A$ выбраны три самые крупные матрицы набора: выполнено 300 повторений на \texttt{add20}, 100 на \texttt{ca-GrQc} и 50 на \texttt{ca-HepTh}. Для $A\times x$ выбраны две самые крупные матрицы: выполнено 1000 повторений на \texttt{ca-GrQc} и 500 на \texttt{ca-HepTh}. Исходные данные и программа профилирования приведены в разделе «Ссылки».

При $A\times A$ просмотр буфера занял 76{,}66\% на \texttt{add20}, 86{,}22\% на \texttt{ca-GrQc} и 90{,}23\% на \texttt{ca-HepTh}. На каждой из трёх матриц эта фаза занимала наибольшую долю времени.

При $A\times x$ фаза двоичного поиска и накопления заняла 93{,}31\% на \texttt{ca-GrQc} и 94{,}63\% на \texttt{ca-HepTh}. За 1000 повторений на \texttt{ca-GrQc} выполнено 28\,980\,000 поисков и 265\,038\,000 сравнений индексов. За 500 повторений на \texttt{ca-HepTh} выполнено 25\,985\,500 поисков и 261\,351\,000 сравнений. В этих измерениях данная фаза занимала наибольшую долю времени.

\section{Заключение}
В ходе работы были получены следующие результаты.
\begin{itemize}
\item Реализованы умножение двух COO-матриц и умножение COO-матрицы на разреженный вектор.
\item Сравнение с CSparse выполнено на девяти матрицах размером от $62\times62$ до $9877\times9877$.
\item CSparse имела меньшее время во всех проведённых измерениях.
\item Профилирование показало, что в $A\times A$ наибольшую долю времени занимает просмотр плотного буфера строки, а в $A\times x$~--- фаза двоичного поиска и накопления.
\end{itemize}

\section*{Ссылки}
\begin{itemize}
\item исходный код и скрипты измерений: \url{https://github.com/rmnsmnsk/implementation-of-linear-algebra-primitives} (дата обращения: 06.09.2026);
\item техническая документация: \url{https://github.com/rmnsmnsk/implementation-of-linear-algebra-primitives/blob/main/README.md} (дата обращения: 06.09.2026);
\item исходные данные профилирования: \url{https://github.com/rmnsmnsk/implementation-of-linear-algebra-primitives/blob/main/profiling_results.csv} (дата обращения: 06.09.2026);
\item необработанный вывод профилировщика: \url{https://github.com/rmnsmnsk/implementation-of-linear-algebra-primitives/blob/main/profiling_raw_wsl.txt} (дата обращения: 06.09.2026);
\item программа профилирования: \url{https://github.com/rmnsmnsk/implementation-of-linear-algebra-primitives/blob/main/src/profile.c} (дата обращения: 06.09.2026);
\item SuiteSparse Matrix Collection: \url{https://sparse.tamu.edu/} (дата обращения: 06.09.2026);
\item репозиторий SuiteSparse и CSparse: \url{https://github.com/DrTimothyAldenDavis/SuiteSparse} (дата обращения: 06.09.2026);
\item репозиторий с отчётами: \url{https://github.com/gsvgit/1_year_practice_2026_reports} (дата обращения: 06.09.2026).
\end{itemize}
\end{document}
